\chapter{Conclusões e Trabalho Futuro}
\label{chap:conclusions}

\begin{introduction}
Este capítulo descreve as conclusões gerais retiradas ao longo do trabalho, bem como sugestões de trabalho futuro.
\end{introduction}


\section{Conclusões}
Este trabalho apresentou o desenvolvimento de uma \textit{pipeline} de processamento em \textit{streaming} de dados de sensores, abordando os desafios de resiliência e integridade, bem como o equilíbrio entre robustez e modularidade.
Em geral, considero que os objetivos propostos foram cumpridos na sua totalidade: consegui desenvolver um sistema baseado em \textit{streaming} capaz de ser estendido para suportar novas fontes de dados com base no caso de uso, verifiquei o seu funcionamento através da recolha de dados de periféricos e \textit{wearables} e considero que, com este sistema, consegui fundamentar a viabilidade do \textit{streaming} num cenário de monitorização contínua envolvendo vários utilizadores durante o uso normal do computador.

\section{Trabalho Futuro}
Apesar dos avanços alcançados, existem ainda funcionalidades e aspetos do trabalho que merecem investigação adicional. Estas incluem: 
\begin{itemize}
    \item Resolver a inferência incorreta de tipos durante a conversão no Kafka Connect (em particular, evitar que \textit{floats} sem casas decimais sejam convertidos em \textit{integers}).
    \item Explorar outras implementações dos protocolos \ac{MQTT} (por exemplo, HiveMQ) e \ac{TSDB} (por exemplo, TimescaleDB) e comparar com as tecnologias adotadas neste trabalho.
    \item Explorar ferramentas de processamento em \textit{streaming} para operações mais complexas como o Kafka Streams.
    \item Suportar a integração de modelos de Inteligência Artificial diretamente no sistema para permitir gerar estimativas mais precisas.
    \item Explorar o protocolo \ac{BLE} para extração de dados de dispositivos \textit{wearables} que não suportem comunicação via \ac{MQTT} (como \textit{smart rings}).
    \item Realizar testes de carga e de escalabilidade para melhor avaliar a viabilidade num cenário real.
    \item Explorar mecanismos de controlo de acesso mais robustos para a publicação e acesso aos dados.
    \item Implementar melhores sistemas de observabilidade que gerem alertas automáticos em caso de falha no sistema.
\end{itemize}

